\documentclass[UTF8,11pt]{ctexart}
\usepackage[a4paper,margin=24mm]{geometry}
\usepackage{amsmath,amssymb,amsthm,mathtools,mathrsfs}
\usepackage[all]{xy}
\usepackage{xurl,hyperref,enumitem}
\usepackage{tikz}
\usetikzlibrary{arrows.meta,cd,decorations.markings,calc,patterns}
\hypersetup{hidelinks,breaklinks=true}
\setlength{\emergencystretch}{3em}
\allowdisplaybreaks
\newtheorem*{theorem}{Theorem}
\newtheorem*{thm}{Theorem}
\newtheorem*{lemma}{Lemma}
\newtheorem*{proposition}{Proposition}
\newtheorem*{prop}{Proposition}
\newtheorem*{corollary}{Corollary}
\newtheorem*{cor}{Corollary}
\newtheorem*{claim}{Claim}
\theoremstyle{definition}
\newtheorem*{definition}{Definition}
\newtheorem*{defn}{Definition}
\newtheorem*{remark}{Remark}
\newtheorem*{example}{Example}
\newcommand{\R}{\mathbb R}
\newcommand{\C}{\mathbb C}
\newcommand{\Z}{\mathbb Z}
\newcommand{\Q}{\mathbb Q}
\newcommand{\N}{\mathbb N}
\newcommand{\F}{\mathbb F}
\newcommand{\D}{\mathbb D}
\newcommand{\bB}{\mathbb B}
\newcommand{\bC}{\mathbb C}
\newcommand{\bR}{\mathbb R}
\newcommand{\bZ}{\mathbb Z}
\newcommand{\bN}{\mathbb N}
\newcommand{\bQ}{\mathbb Q}
\newcommand{\bD}{\mathbb D}
\newcommand{\bF}{\mathbb F}
\newcommand{\CP}{\mathbb{CP}}
\newcommand{\RP}{\mathbb{RP}}
\newcommand{\sO}{\mathscr O}
\newcommand{\sI}{\mathscr I}
\newcommand{\sF}{\mathscr F}
\newcommand{\sG}{\mathscr G}
\newcommand{\sH}{\mathscr H}
\newcommand{\sR}{\mathscr R}
\newcommand{\sM}{\mathscr M}
\newcommand{\sV}{\mathscr V}
\newcommand{\sU}{\mathscr U}
\DeclareMathOperator{\Hom}{Hom}
\DeclareMathOperator{\Ext}{Ext}
\DeclareMathOperator{\Tor}{Tor}
\DeclareMathOperator{\Spec}{Spec}
\DeclareMathOperator{\Proj}{Proj}
\DeclareMathOperator{\supp}{supp}
\DeclareMathOperator{\im}{im}
\DeclareMathOperator{\id}{id}
\DeclareMathOperator{\Tr}{Tr}
\DeclareMathOperator{\tr}{tr}
\DeclareMathOperator{\rank}{rank}
\DeclareMathOperator{\ord}{ord}
\DeclareMathOperator{\dist}{dist}
\DeclareMathOperator{\End}{End}
\DeclareMathOperator{\Aut}{Aut}
\DeclareMathOperator{\Gal}{Gal}
\DeclareMathOperator{\vol}{vol}
\DeclareMathOperator{\Cl}{Cl}
\DeclareMathOperator{\coker}{coker}
\DeclareMathOperator{\codim}{codim}
\DeclareMathOperator{\divergence}{div}
\DeclareMathOperator{\Ric}{Ric}
\DeclareMathOperator{\grad}{grad}
\DeclareMathOperator{\Hess}{Hess}
\newcommand{\abs}[1]{\left\lvert#1\right\rvert}
\newcommand{\sabs}[1]{\lvert#1\rvert}
\newcommand{\babs}[1]{\bigl\lvert#1\bigr\rvert}
\newcommand{\Babs}[1]{\Bigl\lvert#1\Bigr\rvert}
\newcommand{\bbabs}[1]{\biggl\lvert#1\biggr\rvert}
\newcommand{\BBabs}[1]{\Biggl\lvert#1\Biggr\rvert}
\newcommand{\norm}[1]{\left\lVert#1\right\rVert}
\newcommand{\snorm}[1]{\lVert#1\rVert}
\newcommand{\bnorm}[1]{\bigl\lVert#1\bigr\rVert}
\newcommand{\Bnorm}[1]{\Bigl\lVert#1\Bigr\rVert}
\newcommand{\bbnorm}[1]{\biggl\lVert#1\biggr\rVert}
\newcommand{\BBnorm}[1]{\Biggl\lVert#1\Biggr\rVert}
\newcommand{\widebar}[1]{\overline{#1}}
\newcommand{\nicefrac}[2]{\frac{#1}{#2}}
\newcommand{\linnprod}[2]{\langle#1,#2\rangle}
\newcommand{\myindex}[1]{#1}
\newcommand{\glsadd}[1]{}
\newcommand{\avoidbreak}{}
\newcommand{\sourceref}[1]{\textnormal{[原文：\nolinkurl{#1}]}}
\newcommand{\thmref}[1]{\sourceref{#1}}
\newcommand{\lemmaref}[1]{\sourceref{#1}}
\newcommand{\propref}[1]{\sourceref{#1}}
\newcommand{\corref}[1]{\sourceref{#1}}
\newcommand{\defnref}[1]{\sourceref{#1}}
\newcommand{\exampleref}[1]{\sourceref{#1}}
\newcommand{\exerciseref}[1]{\sourceref{#1}}
\newcommand{\figureref}[1]{\sourceref{#1}}
\newcommand{\sectionref}[1]{\sourceref{#1}}
\newcommand{\appendixref}[1]{\sourceref{#1}}
\newcommand{\stacksref}[2]{\href{https://stacks.math.columbia.edu/tag/#1}{#2 [#1]}}


\newcommand{\E}{\mathbb E}
\newcommand{\Prob}{\mathbb P}
\newcommand{\Reff}{\mathcal R_{\mathrm{eff}}}
\newcommand{\eps}{\varepsilon}
\DeclareMathOperator{\diam}{diam}
\DeclareMathOperator{\Var}{Var}
\DeclareMathOperator{\Crit}{Crit}
\newcommand{\dd}{\,\mathrm d}


\begin{document}
\section*{086\quad 关系的隐式唯一确定与Beth显式可定义性}
\subsection*{来源与整理范围}
The Open Logic Project，Model Theory / Interpolation / The Definability Theorem；定义、Beth 定理及两个方向的完整证明。
\par\url{https://github.com/OpenLogicProject/OpenLogic/blob/master/content/model-theory/interpolation/definability.tex}
\par\textbf{恢复方式（整理者说明）：}依据人类原文重新整理以补齐缺失题号；并非旧题证文件的逐字节恢复；2026-09-28。
\subsection*{作者命题与人类证明}

\begin{definition}[Explicit and implicit definability]
Let $\mathcal L$ be a first-order language not containing an $n$-ary predicate
$P$. A set of sentences $\Sigma(P)$ in $\mathcal L\cup\{P\}$ explicitly
defines $P$ if there is an $\mathcal L$-formula $C(x_1,\ldots,x_n)$ such that
\[
 \Sigma(P)\models\forall x_1\cdots\forall x_n
       \bigl(P(x_1,\ldots,x_n)\leftrightarrow C(x_1,\ldots,x_n)\bigr).
\]
Let $P'$ be a fresh predicate of the same arity, and let $\Sigma(P')$ result
from replacing $P$ uniformly by $P'$. The theory implicitly defines $P$ if
\[
 \Sigma(P)\cup\Sigma(P')\models\forall x_1\cdots\forall x_n
       \bigl(P(x_1,\ldots,x_n)\leftrightarrow P'(x_1,\ldots,x_n)\bigr).
\]
Equivalently, for a fixed $\mathcal L$-structure $M$, any two relations
$R,R'\subseteq M^n$ making $(M,R)$ and $(M,R')$ models of the theory must be equal.
\end{definition}
\begin{theorem}[Beth definability]
A set $\Sigma(P)$ of sentences implicitly defines $P$ if and only if it
explicitly defines $P$.
\end{theorem}
\begin{proof}
If $\Sigma(P)$ explicitly defines $P$ by $C$, uniform renaming gives
\[
\begin{aligned}
 \Sigma(P)&\models\forall\mathbf x\,(P(\mathbf x)\leftrightarrow C(\mathbf x)),\\
 \Sigma(P')&\models\forall\mathbf x\,(P'(\mathbf x)\leftrightarrow C(\mathbf x)).
\end{aligned}
\]
The implicit definition follows.

Conversely, assume implicit definability and add fresh constants
$c_1,\ldots,c_n$ to the language. Then
\[
 \Sigma(P)\cup\Sigma(P')\models P(\mathbf c)\to P'(\mathbf c).
\]
By compactness, choose finite sets $\Delta_0\subseteq\Sigma(P)$ and
$\Delta_1\subseteq\Sigma(P')$ with the same entailment. Let $D(P)$ be the
conjunction of every sentence $A(P)$ for which either $A(P)\in\Delta_0$
or its renamed copy $A(P')\in\Delta_1$. Let $D(P')$ be its uniform renamed
copy. Thus
\[
 D(P)\land D(P')\models P(\mathbf c)\to P'(\mathbf c).
\]
Rearrange so that the new predicates occur on opposite sides:
\[
 D(P)\land P(\mathbf c)\models D(P')\to P'(\mathbf c).
\]
Craig interpolation gives a sentence $C(\mathbf c)$ containing neither
$P$ nor $P'$ such that
\[
\begin{aligned}
 D(P)\land P(\mathbf c)&\models C(\mathbf c),\\
 C(\mathbf c)&\models D(P')\to P'(\mathbf c).
\end{aligned}
\]
The first entailment implies
$D(P)\models P(\mathbf c)\to C(\mathbf c)$.
In the second entailment rename $P'$ to $P$. This does not change
$C(\mathbf c)$, which contains neither predicate, and gives
\[
 C(\mathbf c)\models D(P)\to P(\mathbf c).
\]
Therefore $D(P)\models C(\mathbf c)\to P(\mathbf c)$. Combining the two
implications yields
\[
 D(P)\models P(\mathbf c)\leftrightarrow C(\mathbf c).
\]
Since $\Sigma(P)\models D(P)$, monotonicity and generalization over the
fresh constants now give
\[
 \Sigma(P)\models\forall x_1\cdots\forall x_n
        \bigl(P(x_1,\ldots,x_n)\leftrightarrow C(x_1,\ldots,x_n)\bigr).
\]
This is the required explicit definition.
\end{proof}

\subsection*{整理说明与先修范围（非作者正文）}
原生 TeX 的逻辑缩写宏改为标准 LaTeX；用 $\mathbf x,\mathbf c$ 缩写元组，未省略量词或证明方向。作者使用的语义与句子范围按前置定义保持，整理者未增加函数符号版本或无限元关系版本。

先修为一阶逻辑紧致性、符号重命名、全称概括和 Craig 插值定理。与079不同，本题目标是由模型扩张的唯一性提取一个不含待定义谓词的公式；关键工作是复制语言、用紧致性同步有限子理论，再使插值项的两向蕴含返回同一个谓词。079证明插值本身，此题并未把079的辅助引理另计。
\subsection*{可修改的简短标注（非作者正文）}
$c$：一阶结构中关系的可定义性

$a$：模型扩张；语言复制与重命名；紧致性；Craig插值；量词概括

\texttt{cross\_theory\_status = yes}。将“同一底结构的扩张唯一”改写成两个副本理论之间的蕴含，并从共同语言插值项构造显式定义；见有限子理论D及重命名步骤。

全部标注为非穷尽、可修改初稿。
\subsection*{来源许可与编辑归属}
原文由 The Open Logic Project 按 CC BY 4.0 发布；本转录保留作者、来源与编辑归属。许可说明见上游 README.md 和 LICENSE.md。
ChatGPT 加入题号、中文来源说明、整理范围和初稿标注；编辑说明与人类证明分列。
\end{document}
